\section*{Chapter \ref{ch:g5:heat-and-insulation} --- Heat and Insulation}

\begin{solution}{exo:g5:heat-and-insulation:1}
\omterm{def:g5:heat-and-insulation:heat}{Heat} is \omterm{def:g5:everyday-energy:energy}{energy} on the move from a hotter body to a colder one; by
itself it flows only downhill, hot to cold, and stops when the two
\omterm{def:g3:temperature-thermometers:temperature}{temperatures} have met.
\end{solution}

\begin{solution}{exo:g5:heat-and-insulation:2}
``Close the door, you are letting our \omterm{def:g5:heat-and-insulation:heat}{heat} out!'' --- cold does not
flow in; the house's \omterm{def:g5:heat-and-insulation:heat}{heat} flows out toward the colder street.
\end{solution}

\begin{solution}{exo:g5:heat-and-insulation:3}
Fast lanes: the copper pan, the steel spoon. Slow lanes: the wool
scarf, still \omterm{def:g2:air-around-us:air}{air}, the wooden spoon, the duvet.
\end{solution}

\begin{solution}{exo:g5:heat-and-insulation:4}
The bottom must be a fast lane, rushing the stove's \omterm{def:g5:heat-and-insulation:heat}{heat} into the
food; the handle must be a slow lane, keeping that same \omterm{def:g5:heat-and-insulation:heat}{heat} out
of your hand. One tool, both talents, each in its place.
\end{solution}

\begin{solution}{exo:g5:heat-and-insulation:5}
Both stand at the same \omterm{def:g3:temperature-thermometers:temperature}{temperature}, but metal is a fast lane: it
whisks \omterm{def:g5:heat-and-insulation:heat}{heat} out of your warmer hand in a rush, and the rush is the
``icy'' feeling. Skin \omterm{def:g2:measuring-length:measure}{measures} the speed of its own \omterm{def:g5:heat-and-insulation:heat}{heat} loss, not
\omterm{def:g3:temperature-thermometers:temperature}{temperature}.
\end{solution}

\begin{solution}{exo:g5:heat-and-insulation:6}
A fair race changes one thing only --- the wrapping. Identical
mugs, equal starting \omterm{def:g5:heat-and-insulation:heat}{heat} and the same room make sure the scarf
alone explains any difference between the two curves.
\end{solution}

\begin{solution}{exo:g5:heat-and-insulation:7}
After $15$ minutes: the wrapped mug about \qty{41}{\celsius}, the
bare mug about \qty{30}{\celsius}. Both curves head for the room's
\qty{20}{\celsius}: the flow runs until the \omterm{def:g3:temperature-thermometers:temperature}{temperatures} meet, and
insulation only slows the journey.
\end{solution}

\begin{solution}{exo:g5:heat-and-insulation:8}
Puffed feathers trap a thick blanket of still \omterm{def:g2:air-around-us:air}{air} --- the champion
slow lane --- around the bird's warm body. The same card: wool
clothes, feather duvets, double windows (or fleece, insulated
walls).
\end{solution}

\begin{solution}{exo:g5:heat-and-insulation:9}
The thermos neither makes \omterm{def:g5:heat-and-insulation:heat}{heat} nor cold: it closes \omterm{def:g5:heat-and-insulation:heat}{heat}'s lanes ---
emptiness between its two walls, shine to bounce the glow. In
January it slows the chocolate's \omterm{def:g5:heat-and-insulation:heat}{heat} from escaping out; in July it
slows the summer's \omterm{def:g5:heat-and-insulation:heat}{heat} from creeping in. One device, both ways,
because flow is all it manages.
\end{solution}

\begin{solution}{exo:g5:heat-and-insulation:10}
On a hot day, \omterm{def:g5:heat-and-insulation:heat}{heat} tries to flow from the warm \omterm{def:g2:air-around-us:air}{air} \emph{into} the
cold tubs. The wool turns that inward flow into a trickle --- the
same slowing it performs around a warm body. Insulation has no
favorites; the seller, not the customer, understood it.
\end{solution}

\begin{solution}{exo:g5:heat-and-insulation:11}
The double window imprisons a layer of still \omterm{def:g2:air-around-us:air}{air} --- a slow lane
--- between its panes, so the house's \omterm{def:g5:heat-and-insulation:heat}{heat} trickles out where the
single window lets it stream. Less escaping \omterm{def:g5:heat-and-insulation:heat}{heat}, less fuel to
replace it. But insulation only slows the downhill flow, never
stops it: some \omterm{def:g5:heat-and-insulation:heat}{heat} always leaks winter-ward, and the stove must
keep topping it up.
\end{solution}
